\section{Background}


\subsection{Genetic Algorithms}

Genetic algorithms (GAs) \cite{holland1975} are metaheuristic optimization methods inspired by biological evolution. A GA maintains a population of candidate solutions (individuals) and iteratively applies:

\begin{enumerate}
\item \textbf{Selection}: Choose individuals for reproduction based on fitness.
\item \textbf{Crossover}: Combine genetic material from parents to produce offspring.
\item \textbf{Mutation}: Introduce random variations.
\item \textbf{Replacement}: Form next generation from offspring and survivors.
\end{enumerate}

GAs excel at exploring large, discontinuous search spaces where gradient-based methods fail.

\subsection{Multi-Armed Bandits}

Multi-armed bandit (MAB) algorithms \cite{thompson1933} balance exploration and exploitation in sequential decision problems. Applied to marketing, MAB methods allocate traffic to variants based on estimated performance:

\begin{equation}
\text{UCB}_i = \hat{\mu}_i + c\sqrt{\frac{\ln n}{n_i}}
\end{equation}

where $\hat{\mu}_i$ is the estimated conversion rate for variant $i$, $n$ is total trials, and $n_i$ is trials for variant $i$.

MAB methods optimize within a fixed variant set but do not generate new variants.

\subsection{Marketing Campaign Structure}

A digital advertising campaign comprises:

\begin{itemize}
\item \textbf{Creative}: Visual and textual content (images, headlines, body copy).
\item \textbf{Targeting}: Audience definition (demographics, interests, behaviors).
\item \textbf{Placement}: Where ads appear (platform, position).
\item \textbf{Bidding}: How much to pay for impressions or actions.
\end{itemize}

The optimization objective is typically conversion rate (purchases, signups) subject to budget and brand constraints.

